% from ega3-1-fr.tex
\oldpage[III]{133}

de la famille filtrante croissante des sous-$A$-algèbres finies $A_\alpha$
de $A'$, et si chacune des $A_\alpha$ est un anneau local, l'idéal maximal
de $A_\alpha$ est la trace sur $A_\alpha$ de celui de $A_\beta$ pour
$A_\alpha\subset A_\beta$ par le même raisonnement que ci-dessus, donc $A'$
est un anneau local (0, 10.3.1.3).

On notera que si le complété d'un anneau local noethérien $A$ est intègre
(ce qu'on exprime en disant que $A$ est \emph{analytiquement intègre}), $A$
est unibranche. Soient en effet $\mathfrak m$ l'idéal maximal de $A$, $K$
son corps des fractions, $K'$ le corps des fractions de $\widehat A$~; on
a donc
\[
K'=K\otimes_A\widehat A.
\]
Soit $A'_F$ une sous-$A$-algèbre finie de $K$. Le sous-anneau $B_F$ de $K'$
engendré par $\widehat A$ et $A'_F$ est isomorphe à
$A'_F\otimes_A\widehat A$~; c'est un $\widehat A$-module de type fini,
complété de $A'_F$ pour la topologie $\mathfrak m$-adique
(0\textsubscript{I}, 7.3.3 et 7.3.6). Comme $A'_F$ est un anneau semi-local
(Bourbaki, \emph{Alg. comm.}, chap. IV, \S~2, n\textsuperscript{o}~5,
cor. 3 de la prop. 9) et que son complété est intègre, $A'_F$ ne peut avoir
qu'un seul idéal maximal $\mathfrak m'_F$ et on a
$\mathfrak m'_F\cap A=\mathfrak m$~; d'où notre assertion.

\begin{corollary}[4.3.7]
\label{III.4.3.7-fr}
Sous les hypothèses de (\hyperref[III.4.3.5-fr]{4.3.5}), supposons que la
fermeture algébrique de $\mathrm R(Y)$ dans $\mathrm R(X)$ soit de degré
séparable $n$ et que $y\in Y$ soit unibranche. Alors la fibre $f^{-1}(y)$ a
au plus $n$ composantes connexes. En particulier si la fermeture algébrique
de $\mathrm R(Y)$ dans $\mathrm R(X)$ est radicielle sur $\mathrm R(X)$,
$f^{-1}(y)$ est connexe.
\end{corollary}

\begin{proof}
En effet, soit $\mathcal O'_y$ la clôture intégrale de $\mathcal O_y$~; la
fermeture intégrale $\mathcal O''_y$ de $\mathcal O_y$ dans $\mathrm R(X)$
est aussi celle de $\mathcal O'_y$~; mais on sait que si $\mathcal O'_y$ est
un anneau local, $\mathcal O''_y$ est un anneau semi-local dont le nombre
d'idéaux maximaux est au plus égal à $n$ ([13], vol. I, p. 289, th. 22).

Ce corollaire est essentiellement la forme sous laquelle Zariski énonce son
«~théorème de connexion~» pour les schémas algébriques.
\end{proof}

\begin{remark}[4.3.8]
\label{III.4.3.8-fr}
Si on ajoute aux hypothèses de (\hyperref[III.4.3.7-fr]{4.3.7}) l'hypothèse
que $Y$ est normal au point $y$, la fibre $f^{-1}(y)$ est
\emph{géométriquement connexe}, puisque (avec les notations de
(\hyperref[III.4.3.4-fr]{4.3.4})) $g^{-1}(y)$ est réduit à un point $y'$ et
que $k(y')$ est radiciel sur $k(y)$.
\end{remark}

\begin{definition}[4.3.9]
\label{III.4.3.9-fr}
Étant donné un préschéma localement noethérien $Y$, on dit qu'un morphisme de
type fini $f:X\to Y$ est \emph{universellement ouvert} si, pour tout
préschéma irréductible localement noethérien $Y'$, et tout morphisme dominant
$g:Y'\to Y$, toute composante irréductible de $X'=X\times_YY'$ domine $Y'$.
\end{definition}

Si $Y$ est irréductible, cela revient à dire que si $\eta$, $\eta'$ sont les
points génériques de $Y$ et $Y'$ respectivement (de sorte que
$g(\eta')=\eta$), et si on pose $f'=f_{(Y')}$, toute composante irréductible
de $X'$ rencontre $f'^{-1}(\eta')$ (0\textsubscript{I}, 2.1.8)~; cela
implique donc que pour tout ouvert $U$ de $Y$, le morphisme
$f^{-1}(U)\to U$, restriction de $f$, est universellement ouvert.

\begin{corollary}[4.3.10]
\label{III.4.3.10-fr}
Soient $X$, $Y$ deux préschémas localement noethériens intègres.
$f:X\to Y$ un morphisme propre dominant et universellement ouvert. Si la
fermeture algébrique de $\mathrm R(Y)$ dans $\mathrm R(X)$ est radicielle
sur $\mathrm R(Y)$, toute fibre $f^{-1}(y)$ ($y\in Y$) est géométriquement
connexe.
\end{corollary}

\begin{proof}
On peut se borner au cas où $Y=\operatorname{Spec}(B)$, $B$ étant un anneau
intègre noethérien. Il résulte alors de (II, 7.1.7) qu'il existe un anneau
local noethérien intégralement clos $A$ qui domine $\mathcal O_y$ et a
$\mathrm R(Y)$ pour corps des fractions. Soit
$Y'=\operatorname{Spec}(A)$, et soit $h:Y'\to Y$ le morphisme correspondant
à l'injection canonique $B\to A$, qui est birationnel (donc dominant)~; en
outre, si $y'$ est l'unique point fermé de $Y'$, on a $h(y')=y$. Soient
$X'=X\times_YY'$, $f'=f\times1_{Y'}$~; désignons par $\eta$, $\eta'$,
$\xi$ les points génériques de $Y$, $Y'$ et $X$ respectivement, de sorte
que $f(\xi)=\eta$ et $h^{-1}(\eta)=\{\eta'\}$~; en outre,
$k(\eta)=k(\eta')=\mathrm R(Y)$, donc $f'^{-1}(\eta')$ est isomorphe à
$f^{-1}(\eta)$ (I, 3.6.4) et en particulier, puisque $\xi$ est le point
générique de $f^{-1}(\eta)$ (0\textsubscript{I}, 2.1.8),
$f'^{-1}(\eta')$ a un seul point générique. Mais par hypothèse, toute
composante irréductible de $X'$ a son point générique dans
$f'^{-1}(\eta')$, donc $X'$ est nécessairement irréductible, son point
générique $\xi'$ est le point générique de $f'^{-1}(\eta')$ et on a
$k(\xi')=k(\xi)$.

